% Figure: a PCA loadings plot (the rows of V_k) for the
% bridge-monitoring case study. Each of the eight sensors is a point
% in the (PC1, PC2) loadings plane; sensors that move together cluster
% along a common direction, and the two structural groups separate.
\begin{figure}[htbp]
\centering
\begin{tikzpicture}
\begin{axis}[
  width=12cm,
  height=9cm,
  xlabel={loading on PC1 (global thermal-traffic mode)},
  ylabel={loading on PC2 (span-asymmetry mode)},
  xmin=-0.1, xmax=0.6,
  ymin=-0.6, ymax=0.6,
  axis lines=middle,
  enlargelimits=false,
  tick label style={font=\footnotesize},
  label style={font=\small},
  title={PCA loadings: eight strain gauges in the first two principal directions},
  title style={font=\bfseries},
  clip=false,
]
% Unit-circle guide (correlation-loading circle).
\draw[enggray, thin, dashed] (axis cs:0,0) circle [radius=0.5];
% North-span gauges (positive PC2): cluster up-right.
\addplot[only marks, mark=*, mark size=2pt, engblue] coordinates {
  (0.41, 0.34) (0.39, 0.40) (0.44, 0.28) (0.37, 0.45)
};
% South-span gauges (negative PC2): cluster down-right.
\addplot[only marks, mark=square*, mark size=2pt, engred] coordinates {
  (0.40, -0.33) (0.43, -0.27) (0.36, -0.44) (0.42, -0.30)
};
% Loading vectors for two representative gauges.
\draw[engblue, thick, -{Stealth}] (axis cs:0,0) -- (axis cs:0.39,0.40);
\draw[engred, thick, -{Stealth}] (axis cs:0,0) -- (axis cs:0.36,-0.44);
\node[font=\footnotesize, engblue, anchor=south] at (axis cs:0.39,0.42) {N3};
\node[font=\footnotesize, engred, anchor=north] at (axis cs:0.36,-0.46) {S3};
\node[font=\footnotesize, anchor=west] at (axis cs:0.50,0.10) {north span};
\node[font=\footnotesize, anchor=west] at (axis cs:0.50,-0.12) {south span};
\end{axis}
\end{tikzpicture}
\caption[PCA loadings for an eight-gauge bridge-monitoring matrix]{The
loadings plot reads the principal directions back in terms of the
original sensors. Each point is one strain gauge, placed at its loading
$(V_{j1}, V_{j2})$ on the first two principal components. All eight
gauges load positively and almost equally on PC1: the first principal
component is the whole-structure response that every gauge sees
together, the combined thermal expansion and traffic load that lifts
or lowers the entire deck. PC2 splits the gauges into the four north-span
sensors (positive loading) and the four south-span sensors (negative
loading): the second component is the span-asymmetry mode, the part of
the signal that loads one span against the other. The loadings turn an
abstract eigenvector into an engineering statement about which sensors
move together and why, which is what makes PCA a diagnostic rather than
a black box.}
\label{fig:vol02:ch10:pca-loadings}
\end{figure}
